% Source: https://artofproblemsolving.com/community/c3h1351108_geometry_problem
% Stable identity: aops-topic-1351108-q01
Let $BC=8$, $AB=7$, $AC=9$, and drop the altitude $AH=h$ to $BC$; let $BH=x$, so $HC=8-x$. (Since $9^2<7^2+8^2$ and $7^2<8^2+9^2$, the angles at $B$ and $C$ are acute and $H$ lies between them.)

\begin{asy}
usepackage("amsmath");
size(5.5cm);
defaultpen(fontsize(10pt));
path ra(pair P,pair u,pair v,real s=0.2){return P+s*unit(u)--P+s*unit(u)+s*unit(v)--P+s*unit(v);}
pair B=(0,0),C=(8,0),A=(2,3*sqrt(5)),H=(2,0);
draw(A--B--C--cycle,linewidth(0.9)); draw(A--H,red+dashed);
draw(ra(H,C-H,A-H,0.45));
label("$A$",A,N); label("$B$",B,SW); label("$C$",C,SE); label("$H$",H,S);
label("$7$",(A+B)/2,W); label("$9$",(A+C)/2,NE); label("$x$",(B+H)/2,S); label("$8-x$",(H+C)/2,S);
label("$h$",(A+H)/2,E,red);
label("\mbox{}",(min(currentpicture,user=true)+max(currentpicture,user=true))/2,black); // keeps the text after the figure black when LaTeX typesets the labels inline
\end{asy}

By the Pythagorean theorem in the two right triangles,
\[
h^2=49-x^2=81-(8-x)^2=17+16x-x^2 ,
\]
so $16x=32$, $x=2$ and $h=\sqrt{49-4}=\sqrt{45}=3\sqrt5$. The area is
\[
\dfrac{1}{2}\cdot8\cdot3\sqrt5=\boxed{12\sqrt5},
\]
so the student's answer is correct. (Heron's formula agrees: $\sqrt{12\cdot5\cdot4\cdot3}=\sqrt{720}=12\sqrt5$.)
